%% --------------------------------------------------------
\chapter{Математичний вивід тривалості інфляції та вікна спостережень СМВ}
\label{sec:appendix_efolds}
%% --------------------------------------------------------

Для того щоб механізм інфляції розв'язував фундаментальні космологічні проблеми (горизонту та площини), тривалість прискореного розширення повинна перевищувати певне критичне значення. Розглянемо строгий кількісний вивід мінімально необхідної кількості $e$-folds ($N_{\mathrm{total}} \sim 50$--$60$), а також визначимо фізичні межі спостережуваного вікна анізотропії реліктового випромінювання ($N_{\mathrm{CMB}} \sim 7$--$8$).

\section{Вивід загальної тривалості інфляції ($N_{\mathrm{total}}$)}

Проблема горизонту виникає через те, що наш теперішній спостережуваний горизонт (супутня відстань, яку світло пройшло від Великого вибуху до сьогодні $t_0$) є набагато більшим за супутній горизонт Хаббла на початку гарячої стадії Всесвіту ($t_{\mathrm{g}}$).

Умова розв'язання проблеми горизонту вимагає, щоб максимальний супутній розмір причинно-зв'язаної області на самому початку інфляції ($t_{\mathrm{init}}$) був не меншим за супутній радіус нашого теперішнього спостережуваного горизонту:
\begin{equation}
\label{eq:app_cond1}
\frac{1}{a_{\mathrm{init}} H_{\mathrm{inf}}} \ge \frac{1}{a_0 H_0},
\end{equation}
де для спрощення припущено, що параметр Хаббла під час інфляції є приблизно константою ($H_{\mathrm{inf}} \approx \text{const}$).

Перепишемо нерівність \eqref{eq:app_cond1} через відношення масштабних факторів:
\begin{equation}
\label{eq:app_ratio}
\frac{a_{\mathrm{end}}}{a_{\mathrm{init}}} \ge \frac{a_{\mathrm{end}} H_{\mathrm{inf}}}{a_0 H_0}.
\end{equation}
За визначенням, кількість $e$-folds за час інфляції дорівнює $N_{\mathrm{total}} = \ln(a_{\mathrm{end}}/a_{\mathrm{init}})$. Тоді мінімальне значення становить:
\begin{equation}
\label{eq:app_nmin}
  N_{\mathrm{total}} \ge \ln \left( \frac{a_{\mathrm{end}}}{a_0} \right) + \ln \left( \frac{H_{\mathrm{inf}}}{H_0} \right).
\end{equation}

Щоб оцінити відношення $a_{\mathrm{end}}/a_0$, скористаємося законом збереження ентропії у сопутньому об'ємі $S = a^3 s = \text{const}$ від моменту закінчення інфляції та розігріву до сьогоднішнього дня. Густина ентропії релятивістської плазми дорівнює:
\begin{equation}
s = \frac{2\pi^2}{45} g_* T^3,
\end{equation}
де $g_*$ --- ефективна кількість релятивістських ступенів вільності. Нехтуючи зміною $g_*$ за еволюцію (що дає похибку $\Delta N \sim 1$), отримуємо:
\begin{equation}
\label{eq:app_entropy}
\frac{a_{\mathrm{end}}}{a_0} \approx \frac{T_0}{T_{\mathrm{reht}}},
\end{equation}
де $T_0 \approx 2.73~\text{К} \approx 2.3 \times 10^{-4}~\text{еВ}$ --- сучасна температура CMB, а $T_{\mathrm{reht}}$ --- температура розігріву Всесвіту.

Енергетична густина наприкінці інфляції визначається фрідманівським рівнянням $\rho_{\mathrm{inf}} = 3 M_{\mathrm{Pl}}^2 H_{\mathrm{inf}}^2$. Припускаючи швидкий розігрів, коли вся енергія поля переходить у випромінювання, маємо $\rho_{\mathrm{inf}} \approx \rho_{\mathrm{rad}} = \frac{\pi^2}{30} g_* T_{\mathrm{reht}}^4$. Звідси виразимо параметр Хаббла:
\begin{equation}
\label{eq:app_hinf}
H_{\mathrm{inf}} \approx \frac{\pi}{3} \sqrt{\frac{g_*}{10}} \frac{T_{\mathrm{reht}}^2}{M_{\mathrm{Pl}}}.
\end{equation}

Підставимо співвідношення \eqref{eq:app_entropy} та \eqref{eq:app_hinf} у граничне рівняння \eqref{eq:app_nmin}:
\begin{equation}
N_{\mathrm{total}} \ge \ln \left( \frac{T_0}{T_{\mathrm{reht}}} \right) + \ln \left( \frac{\pi}{3} \sqrt{\frac{g_*}{10}} \frac{T_{\mathrm{reht}}^2}{M_{\mathrm{Pl}} H_0} \right).
\end{equation}
Групуючи доданки з $T_{\mathrm{reht}}$, отримуємо фінальну аналітичну формулу для мінімальної тривалості інфляції:
\begin{equation}
\label{eq:app_final_n}
N_{\mathrm{total}} \ge \ln \left( \frac{T_0}{H_0} \right) + \ln \left( \frac{T_{\mathrm{reht}}}{M_{\mathrm{Pl}}} \right) + \mathcal{C},
\end{equation}
де константа $\mathcal{C} = \ln \left( \frac{\pi}{3} \sqrt{\frac{g_*}{10}} \right) \approx 0$ для типових значень $g_* \sim 100$.

Підставимо чисельні значення сучасних спостережуваних констант:
\begin{itemize}
    \item $H_0 \approx 67.4~\text{км/с/Мпк} \approx 1.4 \times 10^{-33}~\text{еВ}$,
    \item $T_0 \approx 2.3 \times 10^{-4}~\text{еВ} \implies \ln(T_0 / H_0) \approx \ln(1.6 \times 10^{29}) \approx 67.2$,
    \item $M_{\mathrm{Pl}} \approx 2.4 \times 10^{18}~\text{ГэВ} = 2.4 \times 10^{27}~\text{еВ}$.
\end{itemize}

Точне значення другого логарифма залежить від невідомого нам масштабу енергії розігріву $T_{\mathrm{reht}}$:
\begin{enumerate}
    \item \textbf{Верхня межа (Масштаб Великого об'єднання):} Якщо інфляція відбувалася на максимальних дозволених енергіях $T_{\mathrm{reht}} \sim 10^{16}~\text{ГэВ}$, то $\ln(T_{\mathrm{reht}}/M_{\mathrm{Pl}}) = \ln(10^{-2}) \approx -4.6$. Тоді:
    \begin{equation}
    N_{\mathrm{total}} \ge 67.2 - 4.6 \approx \mathbf{62.6}.
    \end{equation}
    \item \textbf{Нижня межа (Низькоенергетичний розігрів):} Якщо інфляція відбувалася при значно нижчих енергіях, близьких до елекстрослабкого масштабу або мінімально припустимого за нуклеосинтезом $T_{\mathrm{reht}} \sim 10^{3}~\text{ГэВ}$, то $\ln(T_{\mathrm{reht}}/M_{\mathrm{Pl}}) = \ln(10^{-15}) \approx -34.5$. Тоді:
    \begin{equation}
    N_{\mathrm{total}} \ge 67.2 - 34.5 \approx \mathbf{32.7}.
    \end{equation}
\end{enumerate}
Оскільки більшість реалістичних суперсиметричних та струнних моделей тяжіють до високих масштабів енергій ($10^{13}$--$10^{15}$~ГэВ), у космології як стандартне базове значення приймають вікно **$N_{\mathrm{total}} \sim 50$--$60$**.



\section{Вивід меж вікна спостережень СМВ ($N_{\mathrm{CMB}}$)}

Ми бачимо реліктове випромінювання не з усього простору, а лише всередині нашого поточного горизонту Хаббла. Хвилі первинних квантових збурень інфлатона покидали горизонт під час інфляції, коли їхня фізична довжина хвилі $\lambda_{\mathrm{phys}}(t) = \frac{2\pi}{k} a(t)$ зрівнювалася з радіусом Хаббла $H_{\mathrm{inf}}^{-1}$.

Умова виходу моди з-під горизонту має вигляд:
\begin{equation}
\label{eq:app_exit}
k = a_{\mathrm{exit}} H_{\mathrm{inf}}.
\end{equation}

Найбільший просторовий масштаб, який ми взагалі можемо зафіксувати сьогодні --- це масштаб нашого сучасного горизонту, що відповідає мінімальному хвильовому числу $k_{\mathrm{min}} = a_0 H_0$ (мультиполь анізотропії СМВ $\ell \sim 2$). Мода, що відповідає цьому масштабу, вийшла за горизонт у момент часу $t_{\mathrm{min}}$, коли:
\begin{equation}
a_0 H_0 = a(t_{\mathrm{min}}) H_{\mathrm{inf}}.
\end{equation}

Найменший масштаб (максимальне $k_{\mathrm{max}}$), який надійно зафіксований космічною місією \emph{Planck} в кутовому спектрі температурної анізотропії, обмежений затуханням Сілка на мультиполях $\ell \sim 2500$. Цей масштаб приблизно в $10^3$ разів менший за радіус сучасного горизонту, тобто $k_{\mathrm{max}} \approx 10^3 \, k_{\mathrm{min}} = 10^3 \, a_0 H_0$. Відповідна мода вийшла з-під горизонту в момент $t_{\mathrm{max}}$, коли:
\begin{equation}
10^3 \, a_0 H_0 = a(t_{\mathrm{max}}) H_{\mathrm{inf}}.
\end{equation}

Знайдемо кількість $e$-folds $\Delta N_{\mathrm{CMB}}$, яка минула між моментом виходу найбільшого спостережуваного масштабу ($t_{\mathrm{min}}$) та найменшого ($t_{\mathrm{max}}$):
\begin{equation}
\Delta N_{\mathrm{CMB}} = \ln \left( \frac{a(t_{\mathrm{max}})}{a(t_{\mathrm{min}})} \right) = \ln \left( \frac{10^3 \, a_0 H_0 / H_{\mathrm{inf}}}{a_0 H_0 / H_{\mathrm{inf}}} \right) = \ln(10^3).
\end{equation}
Обчислимо натуральний логарифм:
\begin{equation}
\Delta N_{\mathrm{CMB}} = 3 \cdot \ln(10) \approx 3 \times 2.302 = \mathbf{6.91}.
\end{equation}

Якщо врахувати додаткове розширення вікна за рахунок спектра поляризації та великомасштабної структури (діапазон масштабів розширюється приблизно до $e^8$), математично строгий інтервал усього доступного для людства вікна спостережень первинних збурень становить:
\begin{equation}
N_{\mathrm{CMB}} \sim \mathbf{7 - 8 \text{ \textit{e}-folds}}.
\end{equation}

Усі структури Всесвіту, які ми вивчаємо експериментально, народилися протягом цієї крихітної миті тривалістю менше 8 е-фолдів на самому фіналі інфляції. Решта $40$--$50$ $e$-folds інфляційного потенціалу $V(\phi)$ назавжди стерті для прямих астрофізичних спостережень СМВ, оскільки відповідні їм квантові флуктуації розтягнуті до масштабів, що лежать глибоко за межами сучасного метагалактичного горизонту.
